%% ma: L10-B13-0005
%% lop: 10
%% chuong: 5
%% bai: 13
%% dang: 10.13.9
%% loai: TL
%% muc: VDC
%% dap_an: "-"
%% nguon: "chuyen-de-cac-so-dac-trung-cua-mau-so-lieu-khong-ghep-nhom-toan-10-chuong-trinh-moi.pdf (GV. Phan Nhật Linh) (P02), Bài 02 Dạng 1 Bài tập tự luận 5"
%% kien_thuc: "Lớp 10 Bài 13 (số trung bình có trọng số), bất đẳng thức x^2+y^2+z^2 >= (x+y+z)^2/3 (chứng minh bằng tổng bình phương)"
%% trang_thai: chinh-thuc
%% ghi_chu: "Đề gốc hỏi 'lợi nhuận bình quân' nhưng bảng chỉ có giá bán; đổi thành giá bán trung bình (doanh thu bình quân 1 kg), bổ sung x, y, z >= 0; đáp số 314 khớp gốc. Tự duyệt theo yêu cầu giáo viên 22/10/2026."
\begin{ex}
Một cửa hàng bán ba loại hoa quả. Trong một quý, sau khi giảm giá mỗi loại lần lượt $x$, $y$, $z$ (nghìn đồng/kg) với $x,y,z\ge0$ và $x+y+z=90$, số liệu bán hàng như sau:
\begin{center}
\begin{tabular}{|c|c|c|c|}\hline
 & Táo & Dưa & Bưởi\\\hline
Giá bán 1 kg (nghìn đồng) & $250-x$ & $350-y$ & $400-z$\\\hline
Số lượng bán (kg) & $250+x$ & $350+y$ & $400+z$\\\hline
\end{tabular}
\end{center}
Tính giá bán trung bình của $1$ kg hoa quả lớn nhất có thể đạt được trong quý (làm tròn đến hàng đơn vị, đơn vị nghìn đồng).
\loigiai{Tổng khối lượng bán: $250+x+350+y+400+z=1090$ (kg), không đổi.

Tổng doanh thu $P=(250^2-x^2)+(350^2-y^2)+(400^2-z^2)=345\,000-\left(x^2+y^2+z^2\right)$.

Vì $3\left(x^2+y^2+z^2\right)-(x+y+z)^2=(x-y)^2+(y-z)^2+(z-x)^2\ge0$ nên $x^2+y^2+z^2\ge\dfrac{90^2}3=2700$, dấu bằng khi $x=y=z=30$.

Suy ra $P\le342\,300$, giá bán trung bình $\dfrac{P}{1090}\le\dfrac{342\,300}{1090}\approx314$ (nghìn đồng/kg).}
\end{ex}
